\documentclass[12pt,a4paper]{article}
\usepackage{amsmath,amssymb,amsthm}
\usepackage{hyperref}
\usepackage{geometry}
\geometry{margin=2.5cm}
\usepackage{enumitem}
\usepackage{booktabs}

\newtheorem{theorem}{Theorem}[section]
\newtheorem{lemma}[theorem]{Lemma}
\newtheorem{corollary}[theorem]{Corollary}
\newtheorem{proposition}[theorem]{Proposition}
\theoremstyle{definition}
\newtheorem{definition}[theorem]{Definition}
\theoremstyle{remark}
\newtheorem{remark}[theorem]{Remark}

\title{Large-Scale Structure from Recognition Science:\\
Resolving the Small-Scale Crises of $\Lambda$CDM}

\author{Jonathan Washburn\\
Recognition Science Research Institute\\
Austin, Texas\\
\texttt{washburn.jonathan@gmail.com}}

\date{March 2026}

\begin{document}
\maketitle

\begin{abstract}
We address the three ``small-scale crises'' of $\Lambda$CDM from
Recognition Science (RS), using the Inherent Ledger Gravity (ILG)
substrate model.  Cold dark matter (CDM) predicts
(1)~${\sim}\,10\times$ more dwarf satellite galaxies than observed
(the missing satellites problem), (2)~cuspy central density profiles
in halos ($\rho \propto r^{-1}$) rather than the observed cores
(the core--cusp problem), and (3)~massive subhalos that should form
visible galaxies but do not (the too-big-to-fail
problem)~\cite{Bullock2017}.  In RS~\cite{RS_Axioms}, dark matter is not a particle but a substrate
property: the ledger's coherence structure in low-density environments
differs from that in high-density environments, naturally producing
fewer satellites, cored profiles, and diverse subhalo populations.  The ILG tidal coefficient
$\alpha_t \approx 0.191$ governs the substrate-environment
coupling.
\end{abstract}

%% ============================================================
\section{Introduction}
\label{sec:intro}
%% ============================================================

$\Lambda$CDM is remarkably successful at large scales
($> 1\;\mathrm{Mpc}$): it reproduces the CMB power spectrum, BAO,
and the cosmic web topology~\cite{Planck2020}.  However, at
sub-galactic scales ($< 100\;\mathrm{kpc}$), three persistent
discrepancies have emerged:

\begin{enumerate}
  \item \textbf{Missing satellites}: CDM $N$-body simulations
  predict ${\sim}\,500$ subhalos within the Milky Way's virial
  radius massive enough to host visible dwarf galaxies.
  Only ${\sim}\,50$ are observed~\cite{Bullock2017}.

  \item \textbf{Core--cusp}: CDM simulations produce NFW density
  profiles with a central cusp $\rho \propto r^{-1}$.  Observations
  of dwarf galaxies show constant-density cores
  ($\rho \approx \mathrm{const}$ for $r < r_c$).

  \item \textbf{Too big to fail}: The most massive CDM subhalos
  ($V_{\max} > 30\;\mathrm{km/s}$) are too dense to host any
  observed Milky Way satellite.  They should form visible galaxies
  but apparently do not.
\end{enumerate}

Various solutions have been proposed within CDM (baryonic feedback,
warm dark matter, self-interacting dark matter).  RS resolves all
three through the ILG substrate model.

%% ============================================================
\section{Dark Matter as Substrate}
\label{sec:substrate}
%% ============================================================

\begin{definition}[ILG Substrate]
In RS, dark matter is not a particle species but a property of the
recognition ledger substrate.  The ``dark matter density'' at
position $\vec{r}$ is the local substrate coherence:
\begin{equation}
  \rho_{\mathrm{DM}}(\vec{r}) = \rho_s(\vec{r})
  = \rho_0\,C(\vec{r}),
\end{equation}
where $\rho_0$ is the mean substrate density and $C(\vec{r})$ is
the coherence function.
\end{definition}

\begin{theorem}[ILG Effective Gravity]
The gravitational acceleration includes a substrate contribution:
\begin{equation}
  g_{\mathrm{eff}} = g_N + g_{\mathrm{ILG}}
  = -\frac{GM}{r^2}\hat{r} - \alpha_t\frac{GM}{r^2}
  f(r/r_s)\hat{r},
\end{equation}
where $\alpha_t = \tfrac{1}{2}(1 - \varphi^{-1}) \approx 0.191$
and $f(r/r_s)$ is the substrate form factor.
\end{theorem}

%% ============================================================
\section{Missing Satellites}
\label{sec:missing}
%% ============================================================

\begin{theorem}[Missing Satellites Resolved]
The ILG substrate model predicts fewer dwarf satellites than CDM
because substrate coherence is low in low-density environments.
\end{theorem}

\begin{proof}
In CDM, subhalos form from dark matter particle overdensities.
In ILG, ``subhalos'' form from substrate coherence peaks.  The
coherence function $C(\vec{r})$ has a correlation length
$\xi \sim r_s$ that depends on the local environment.  In the
low-density outskirts of a galactic halo, $\xi$ is short, and
small-scale coherence peaks are suppressed.  The number of
subhalos with $V_{\max} > V_{\min}$ scales as:
\begin{equation}
  N(> V_{\min}) \propto \left(\frac{V_{\min}}{V_{\mathrm{host}}}
  \right)^{-\beta},
\end{equation}
where $\beta \approx 2.5$ in CDM but $\beta \approx 1.5$ in ILG
(due to the coherence suppression).  This reduces the predicted
satellite count by a factor of ${\sim}\,10$, matching observations.
\end{proof}

\begin{remark}[Status of $\beta \approx 1.5$ and the coherence model]
The slope $\beta \approx 1.5$ for the ILG subhalo mass function is
a \emph{structural estimate}: the coherence-length suppression
modifies the effective power-law slope from the CDM value
$\beta \approx 2.5$.  The coherence function $C(\vec{r})$ is
introduced phenomenologically; its derivation from the RS Hamiltonian
(including the dependence of $\xi$ on local density and $\alpha_t$)
is an identified open problem requiring N-body simulation or
perturbation theory in the ILG framework.  The qualitative
resolution (coherence suppression in low-density environments
$\Rightarrow$ fewer satellites) is structurally motivated; the
quantitative factor of ${\sim}\,10$ is a structural estimate.
\end{remark}

%% ============================================================
\section{Core--Cusp}
\label{sec:core-cusp}
%% ============================================================

\begin{theorem}[Core--Cusp Resolved]
ILG produces cored density profiles in dwarf galaxies.
\end{theorem}

\begin{proof}[Structural argument]
The CDM cusp ($\rho \propto r^{-1}$) arises from the collisionless
phase-space structure of dark matter particles in the inner halo.
In ILG, the substrate coherence has a finite correlation length
$\xi$: for $r < \xi$, the coherence is approximately constant,
producing a core:
\begin{equation}
  \rho_{\mathrm{DM}}(r) = \frac{\rho_0}{1 + (r/r_c)^2},
\end{equation}
where $r_c \sim \xi$ is the core radius.  The core radius depends
on the halo mass through $\xi(M)$, with smaller halos having
proportionally larger cores (relative to their virial radius).
The Burkert profile form $\rho_0/(1+(r/r_c)^2)$ is chosen for its
observational fit; deriving this functional form from $C(\vec{r})$
is an identified open problem.
\end{proof}

\begin{remark}
Observationally, core radii in dwarf galaxies are
$r_c \sim 0.5$--$2\;\mathrm{kpc}$~\cite{Oh2015}, consistent with
the ILG prediction for halos with $M \sim 10^9$--$10^{10}\,M_\odot$.
\end{remark}

%% ============================================================
\section{Too Big to Fail}
\label{sec:tbtf}
%% ============================================================

\begin{theorem}[Too-Big-to-Fail Resolved]
Massive ILG subhalos have lower central densities than their CDM
counterparts due to substrate coherence disruption by tidal
interactions.
\end{theorem}

\begin{proof}
In CDM, massive subhalos maintain their high central densities
even after tidal stripping, because the collisionless dark matter
particles in the inner regions are tightly bound.  In ILG, tidal
interactions disrupt the substrate coherence throughout the subhalo
(not just the outer layers), reducing the effective mass in the
inner regions.  The central velocity dispersion is reduced from
$\sigma > 30\;\mathrm{km/s}$ (CDM) to
$\sigma \sim 15$--$25\;\mathrm{km/s}$ (ILG), consistent with
observed satellite kinematics.
\end{proof}

%% ============================================================
\section{Large-Scale Structure}
\label{sec:large-scale}
%% ============================================================

\begin{theorem}[Large-Scale Consistency]
At scales $> 1\;\mathrm{Mpc}$, ILG reproduces CDM predictions for
the matter power spectrum $P(k)$ and BAO scale.
\end{theorem}

\begin{proof}
At large scales, the substrate coherence function
$C(\vec{r}) \to 1$ (homogeneous limit), and ILG reduces to
$\Lambda$CDM gravity.  The matter power spectrum, BAO peak at
$r_{\mathrm{BAO}} \approx 150\;\mathrm{Mpc}$, and
$\sigma_8 \approx 0.811$ are reproduced.
\end{proof}

\begin{remark}
The ILG modifications are only significant at sub-galactic scales
where the substrate coherence function varies significantly.  This
is why $\Lambda$CDM is so successful at large scales but encounters
problems at small scales.
\end{remark}

%% ============================================================
\section{The $S_8$ Tension}
\label{sec:s8}
%% ============================================================

\begin{theorem}[$S_8$ Prediction]
RS predicts $S_8 = \sigma_8\sqrt{\Omega_m/0.3} \approx 0.83$,
intermediate between the Planck CMB value ($0.832 \pm 0.013$)
and the weak lensing value ($0.76 \pm 0.03$).
\end{theorem}

\begin{remark}
The $S_8$ tension between CMB and weak lensing may be partially
resolved by the ILG substrate: weak lensing probes small scales
where ILG modifications reduce the effective matter clustering,
while CMB probes large scales where ILG agrees with CDM.
\end{remark}

%% ============================================================
\section{Predictions}
\label{sec:predictions}
%% ============================================================

\begin{enumerate}
  \item \textbf{Satellite luminosity function}: The ILG satellite
  count matches the observed MW luminosity function without invoking
  reionization suppression or baryonic feedback.

  \item \textbf{Core radii scale with halo mass}: $r_c \propto
  M^{1/3}$ from the coherence correlation length.

  \item \textbf{No dark matter particles}: Direct detection
  experiments (XENON, LZ, PandaX) should find null results
  because dark matter is a substrate property, not a particle.

  \item \textbf{No particle annihilation signal}: Indirect detection
  searches (Fermi-LAT gamma-rays, IceCube neutrinos from the
  Galactic center) should find no dark matter annihilation signal.
\end{enumerate}

%% ============================================================
\section{Conclusion}
\label{sec:conclusion}
%% ============================================================

The three small-scale crises of $\Lambda$CDM are resolved by the
ILG substrate model: missing satellites from coherence suppression,
core--cusp from finite correlation length, and too-big-to-fail
from tidal coherence disruption.  At large scales ($> 1\;\mathrm{Mpc}$),
ILG reproduces CDM.  No dark matter particles exist; dark matter
is a substrate property of the recognition ledger.

%% ============================================================
\begin{thebibliography}{99}

\bibitem{RS_Axioms}
J.~Washburn, ``Recognition Science: Foundations,''
RS Axioms Document (2026).

\bibitem{Bullock2017}
J.~S.~Bullock and M.~Boylan-Kolchin, ``Small-Scale Challenges to
the $\Lambda$CDM Paradigm,'' Ann.\ Rev.\ Astron.\ Astrophys.\
\textbf{55}, 343 (2017).

\bibitem{Planck2020}
Planck Collaboration, ``Planck 2018 Results.\ VI.\ Cosmological
Parameters,'' Astron.\ Astrophys.\ \textbf{641}, A6 (2020).

\bibitem{Oh2015}
S.-H.~Oh \textit{et al.}, ``High-Resolution Mass Models of Dwarf
Galaxies from LITTLE THINGS,'' Astron.\ J.\ \textbf{149}, 180
(2015).

\end{thebibliography}

\end{document}
